\chapter{Introduction}

\section{Research Background}

Computed tomography (CT) is one of the most widely used medical imaging modalities because it provides high-resolution cross-sectional information about internal anatomical structures. It is used in diagnosis, screening, treatment planning, and follow-up examinations for many clinical tasks. However, CT imaging relies on X-ray radiation. Although the radiation dose of a single scan is usually controlled within clinical guidelines, cumulative exposure becomes a concern for repeated follow-up, pediatric imaging, population screening, and patients who require long-term monitoring. Reducing CT radiation dose while preserving image quality is therefore an important problem in medical imaging.

Low-dose CT (LDCT) imaging reduces the radiation burden by lowering the tube current, exposure time, or related acquisition settings. The cost of dose reduction is a higher noise level in the measured projection data and, after reconstruction, stronger noise and artifacts in the image domain. LDCT images often contain granular noise, streak-like patterns, and local texture degradation. These degradations can reduce the visibility of low-contrast lesions, make organ boundaries less reliable, and affect the confidence of image interpretation. Therefore, LDCT denoising is not merely a cosmetic image enhancement problem. It is a restoration task in which anatomical structure must be preserved while noise is suppressed.

Traditional denoising and reconstruction approaches include filtering, statistical iterative reconstruction, dictionary learning, and model-based regularization. These methods can be effective, but they often require careful parameter tuning and may struggle to balance noise removal with detail preservation. Deep learning has changed the landscape of medical image restoration. Convolutional neural networks (CNNs), such as RED-CNN, learn an image-domain mapping from LDCT to normal-dose CT (NDCT) and have become strong baselines for LDCT denoising \cite{redcnn}. More recently, Transformer-based models have been introduced into medical image denoising and general image restoration because attention mechanisms can model longer-range dependencies than local convolution alone \cite{attention,ctformer,restormer}.

However, higher PSNR or SSIM does not automatically mean that a denoised CT image is clinically safer. A model may improve global pixel metrics by producing a smoother image, but excessive smoothing can remove small structures or weaken edges. This is especially important for medical imaging, where the removed signal should be noise-like rather than anatomy-like. Motivated by this issue, this thesis evaluates denoising models not only through PSNR, SSIM, and RMSE, but also through residual noise maps. The central question becomes: what does the model remove from the LDCT input?

\section{Research Objectives}

This thesis studies supervised LDCT image denoising with paired LDCT and NDCT slices. The project has four main objectives. First, it aims to build a complete and reproducible LDCT denoising pipeline, including data preparation, patient-level splitting, model training, checkpoint management, full-slice inference, metric computation, and visualization. Second, it adapts a Restormer-style image restoration model to LDCT denoising and refers to this adapted model as CTRestormer. Third, it compares CTRestormer with a CNN baseline and a Transformer-based CT denoising baseline under the same evaluation setting. Fourth, it studies residual maps to interpret whether each model removes mainly dose-induced noise or also removes structure-related signal.

The task is formulated as paired image restoration. Given a low-dose CT slice \(x\), a model \(f_{\theta}\) predicts a denoised output \(\hat{y}=f_{\theta}(x)\), which should approximate the normal-dose target \(y\). The main performance metrics are PSNR, SSIM, and RMSE. In addition, this thesis defines the estimated LDCT residual as \(x-y\), the removed residual as \(x-\hat{y}\), and the prediction error as \(\hat{y}-y\). These residual quantities provide a more detailed view of denoising behavior.

\section{Main Contributions}

The main contributions of this thesis are summarized as four parts.

First, this thesis constructs an end-to-end LDCT denoising workflow. The workflow starts from paired AAPM-Mayo LDCT and NDCT slices, uses a patient-level test split, trains multiple denoising models, performs full-resolution inference, computes quantitative metrics, and generates figures and residual maps. This makes the project more than an isolated model run; it is a complete experimental system.

Second, this thesis adapts a Restormer-style restoration Transformer to the LDCT denoising task. The proposed CTRestormer uses restoration-oriented Transformer blocks, hierarchical feature processing, residual prediction, patch-based training, and tiled inference. These design choices are important because full \(512 \times 512\) CT slices are memory intensive for Transformer models. The implementation records practical strategies for running the model under limited local hardware.

Third, this thesis provides a fair comparison across models. RED-CNN is used as a representative CNN-based LDCT denoising baseline, while CTformer is used as a Transformer-based CT denoising baseline. CTRestormer is evaluated at several checkpoints to show training progression instead of relying on a single isolated result.

Fourth, this thesis introduces residual noise analysis as an additional interpretation layer. Instead of only looking at denoised images or aggregate metrics, the thesis compares what each model removes from the input. Residual correlation, edge leakage, histogram shape, and power spectrum are used to examine whether the removed residual resembles estimated LDCT noise or contains anatomical edges. This analysis provides evidence about possible over-smoothing risk, while still acknowledging that residual maps are not a substitute for clinical evaluation.

\section{Thesis Organization}

The remainder of the thesis is organized as follows. Chapter 2 reviews related work on LDCT denoising, CNN-based restoration, Transformer-based image restoration, and self-supervised denoising. Chapter 3 presents the methodology, including problem formulation, dataset preparation, baseline models, CTRestormer, and residual noise analysis. Chapter 4 describes the experimental setup, including the dataset split, training configuration, inference protocol, and evaluation metrics. Chapter 5 reports quantitative and qualitative results. Chapter 6 discusses residual noise analysis in detail and interprets model behavior beyond standard metrics. Chapter 7 concludes the thesis and outlines future work.
